\documentclass[a4paper,14pt]{extarticle}

\usepackage[T2A]{fontenc}
\usepackage[utf8]{inputenc}
\usepackage[russian]{babel}
\usepackage{geometry}
\geometry{left=3cm,right=1.5cm,top=2cm,bottom=2cm}
\usepackage{setspace}
\onehalfspacing
\usepackage{indentfirst}
\usepackage{mathptmx}
\pagestyle{empty}
\setlength{\parindent}{1.25cm}
\setlength{\parskip}{0pt}

\begin{document}

\begin{center}
\textbf{%
Отзыв научного руководителя\\[0.5\baselineskip]
на диссертацию Горбова Григория Владиславовича\\[0.5\baselineskip]
«Методы и алгоритмы безопасного обучения с подкреплением в робототехнических задачах»,\\[0.5\baselineskip]
представленную на соискание учёной степени кандидата технических наук\\
по специальности 1.2.1~--- «Искусственный интеллект и машинное обучение».%
}
\end{center}

Разработка методов безопасного обучения с подкреплением является актуальным направлением в области искусственного интеллекта и робототехники. Современные автономные агенты должны не только достигать поставленных целей, но и соблюдать ограничения безопасности при взаимодействии со сложной и частично неизвестной средой. Эта задача особенно важна для робототехнических систем, поскольку ошибки в процессе обучения и исследования среды могут приводить к столкновениям, повреждению объектов или нарушению правил, заданных пользователем.

Существующие методы безопасного обучения с подкреплением обладают рядом ограничений. Одноуровневые безопасные стратегии часто становятся чрезмерно консервативными и перестают эффективно исследовать среду. Методы на основе моделей мира позволяют заранее оценивать последствия действий, однако требуют надёжной оценки будущих штрафов и долгосрочного риска. Иерархические методы помогают решать задачи с большим горизонтом, но сами по себе не гарантируют безопасность выбираемых подцелей и траекторий их достижения. Кроме того, при переходе к более реалистичным сценариям ограничения безопасности могут задаваться не только численно, но и на естественном языке. Всё это определяет актуальность диссертационной работы Григория.

В диссертационной работе последовательно рассматриваются и решаются указанные проблемы. Григорием проведён анализ современных методов безопасного обучения с подкреплением, включая методы оптимизации стратегии с ограничениями, подходы на основе моделей мира, безопасное иерархическое обучение и задачи с естественно-языковыми ограничениями. На основе этого анализа сформулированы основные направления исследования: оценка будущих накопленных штрафов нарушений, улучшение безопасного исследования за счёт подцелей, обеспечение безопасности в иерархической структуре и построение бенчмарка для задач с языковыми ограничениями.

Одним из результатов работы является исследование компонентов модели мира, отвечающих за оценку штрафов и долгосрочного риска. В диссертации показано, что для безопасного планирования недостаточно учитывать только мгновенный штраф на коротком горизонте. Необходима оценка будущих накопленных штрафов, а также устойчивый способ агрегации таких оценок. Полученные результаты показывают, что корректная оценка долгосрочного риска позволяет улучшить баланс между достижением цели и соблюдением ограничений безопасности.

Также в работе предложен метод SG-Safe, направленный на улучшение безопасного исследования в задачах с большим горизонтом. Метод использует генерацию промежуточных подцелей как дополнительный обучающий сигнал для безопасной стратегии. Такой подход позволяет агенту активнее исследовать среду и находить траектории к цели, не сводя поведение к чрезмерно осторожному избеганию любых потенциально опасных действий.

Дальнейшим развитием этого направления стал метод ITES, в котором безопасность учитывается как на уровне выбора подцелей, так и на уровне их достижения. Для этого используются модель штрафов и модель мира, позволяющие предварительно оценивать риск промежуточных целей и действий низкоуровневой стратегии. Предложенный метод позволяет решать робототехнические задачи с большим горизонтом, где агенту необходимо не только разбивать поведение на промежуточные этапы, но и контролировать безопасность каждого такого этапа.

Отдельным результатом диссертационной работы является бенчмарк Caged CrafText, предназначенный для оценки безопасного обучения с подкреплением в задачах, где цель и ограничения безопасности задаются на естественном языке. Бенчмарк включает различные типы ограничений и позволяет анализировать способность агента понимать инструкцию, связывать её с объектами и состояниями среды, планировать поведение и соблюдать требования безопасности. Этот результат расширяет класс задач безопасного обучения с подкреплением и создаёт основу для дальнейшего исследования агентов, работающих с языковыми правилами поведения.

Представленные в диссертации результаты обладают достоверностью за счёт проведения экспериментальной проверки на нескольких типах сред, сравнения с современными базовыми методами и выполнения абляционного анализа. Работа имеет как научную, так и практическую значимость, поскольку предложенные методы могут быть использованы при разработке автономных робототехнических систем, которым необходимо обучаться и адаптироваться к новым условиям при ограничении числа небезопасных действий.

Основные результаты диссертационной работы опубликованы в рецензируемых научных изданиях и материалах международной конференции. В частности, результаты представлены в работах по методам SG-Safe и SPOWL, а также по бенчмарку Caged CrafText. Результаты по методу ITES оформлены в виде рукописи и находятся на стадии рецензирования. Публикации подтверждают научную новизну и значимость полученных результатов.

В ходе диссертационного исследования Григорий показал высокий уровень самостоятельности и профессиональной подготовки. Он продемонстрировал способность формулировать научные задачи, разрабатывать алгоритмические решения, реализовывать экспериментальные системы и анализировать полученные результаты. Работа отличается системностью и охватывает несколько связанных направлений: безопасное обучение с подкреплением, модели мира, иерархическое обучение и мультимодальные задачи с естественно-языковыми ограничениями.

Диссертация Горбова Григория Владиславовича «Методы и алгоритмы безопасного обучения с подкреплением в робототехнических задачах» является завершённым научным исследованием, обладающим научной новизной, теоретической значимостью и практической ценностью. Считаю, что данная диссертационная работа удовлетворяет всем требованиям Положения МФТИ о порядке присуждения учёных степеней, а её автор, Горбов Григорий Владиславович, заслуживает присуждения учёной степени кандидата технических наук по специальности 1.2.1~--- «Искусственный интеллект и машинное обучение».

\vspace{1\baselineskip}
\noindent Научный руководитель,\\
доктор физико-математических наук, профессор\\
Панов Александр Игоревич

\vspace{2\baselineskip}
\noindent «\underline{\hspace{1cm}}» \underline{\hspace{3cm}} 2026~г.

\end{document}
